
Una esterilización se considera efectiva cuando es suficiente para destruir las esporas de \sci{Geobacillus stearothermophilus}, que produce las esporas más resistentes conocidas hasta el momento.

\begin{figure}\centering
    \resizebox{\linewidth}{!}{\def\xlong{20}
\def\ylong{10}
\def\Tmax{0.9*\ylong}
\def\ti{0.35*\xlong}
\def\tf{0.65*\xlong}

\begin{tikzpicture}[
    texto/.style={black, font=\sffamily},
    lim ancho/.style={text width=2.5cm, align=center},
    eje/.style={thick, arrows=-Stealth},
    limites/.style={dashed, DarkRed},
    marca/.style={loosely dotted, DarkRed},
    linea central/.style={ultra thick, DarkBlue}
]

\draw[eje] (0,0) -- (\xlong,0) node[midway, below=5pt, texto] {Tiempo (min)};
\draw[eje] (0,0) -- (0,\ylong) node[pos=0.6, rotate=90, left=10pt, texto] {Temperatura (\unit{\kelvin})};

\draw[linea central] (0,1) to[out=60,in=200] (\ti,\Tmax) -- (\tf,\Tmax) to[out=-60,in=160] (\xlong,1);

\draw[limites] (\ti,0) -- (\ti,1.1*\Tmax);
\draw[limites] (\tf,0) -- (\tf,1.1*\Tmax);

\draw[marca] (0,\Tmax) -- (\ti,\Tmax) node[midway, above, texto, lim ancho] {Temperatura de sostenimiento};

\node[texto, lim ancho] at (0.6*\ti, 0.4*\Tmax) {Fase de calentamiento};
\node[texto, lim ancho] at (0.5*\xlong, 0.6*\Tmax) {Fase de sostenimiento};
\node[texto, lim ancho] at (0.5*\xlong+0.6*\tf, 0.6*\Tmax) {Fase de enfriamiento};

\end{tikzpicture}}
    \rule{\linewidth}{0.5pt}
    \caption{Esquema habitual de temperaturas en un proceso de esterilización}
    \label{fig:temperaturas_de_esterilizacion}
\end{figure}

\begin{equation}
    k = A \cdot e^{-\frac{E_{a}}{R \cdot T}}
\end{equation}

Donde $k$ es la constante de Arrhenius, $A$ es la constante preexponencial, $E_{a}$ es la energía de activación, o inactivación en casos de esterilización, $R$ es la constante de los gases ideales y $T$ es la temperatura en Kelvin (\unit{\kelvin}).

A partir de esto se calcula $\nabla$

\begin{equation}
    \nabla = k \cdot (t_{1} - t_{0}) = k \cdot \Delta T 
\end{equation}

\begin{equation}
    \nabla = \ln(\frac{N_{0}}{N})
\end{equation}

Donde $N_{0}$ son las esporas iniciales y $N$ la cantidad de esporas con que se desea terminar.

\begin{equation}
    \nabla_{\text{Total}} = \nabla_{\text{Calentamiento}} + \nabla_{\text{Sostenimiento}} + \nabla_{\text{Enfriamiento}}
\end{equation}

\begin{figure}\centering
    \resizebox{\linewidth}{!}{\def\xlong{20}
\def\ylong{10}
\def\kmax{0.9*\ylong}
\def\ti{0.3*\xlong}
\def\tf{0.65*\xlong}

\begin{tikzpicture}[
    texto/.style={black, font=\sffamily},
    lim ancho/.style={text width=2.5cm, align=center},
    eje/.style={thick, arrows=-Stealth},
    limites/.style={dashed, DarkRed},
    marca/.style={loosely dotted, DarkRed},
    linea central/.style={ultra thick, DarkBlue}
]

\draw[eje] (0,0) -- (\xlong,0) node[midway, below, texto] {Tiempo ($t$)};
\draw[eje] (0,0) -- (0,\ylong) node[midway, rotate=90, above, texto] {Constante de velocidad ($k$)};

\draw[linea central] (0,0) to [out=0,in=-90] (\ti,\kmax) -- (\tf,\kmax) to[out=-80,in=140] (\xlong,0);

\draw[limites] (\ti,0) -- (\ti,1.1*\kmax);
\draw[limites] (\tf,0) -- (\tf,1.1*\kmax);

\node[texto, lim ancho] at (0.4*\ti, 0.5*\kmax) {$\nabla_{\text{Calentamiento}}$};
\node[texto, lim ancho] at (0.5*\xlong, 0.6*\kmax) {$\nabla_{\text{Sostenimiento}}$};
\node[texto, lim ancho] at (\tf+0.4*\tf, 0.45*\kmax) {$\nabla_{\text{Enfriamiento}}$};

\end{tikzpicture}}
    \rule{\linewidth}{0.5pt}
    \caption{Esquema habitual de temperaturas en un proceso de esterilización}
    \label{fig:constante_de_arrhenius_en_esterilizacion}
\end{figure}